\begin{center}\begin{tikzpicture}[font=\small,>=stealth,line join=round]
\draw[gray!40,rounded corners=3pt] (0,0) rectangle (14.2,8.3);\node at (7.1,7.700000000000001) {例5：同是奇函数，零处是否有定义却不同};
\begin{scope}[shift={(4.1,3.9499999999999997)},x=1.0714285714285714cm,y=0.575cm]\draw[->,gray] (-2.8,0)--(2.8,0) node[right]{$x$};\draw[->,gray] (0,-4)--(0,4) node[above]{$y$};\draw[blue,thick,solid,domain=-2.5:2.5,samples=180] plot(\x,{\x/(1+(\x)^2)});\draw[violet,thick,dashed,domain=-2.5:-0.35,samples=180] plot(\x,{\x+1/\x});\draw[violet,thick,dashed,domain=0.35:2.5,samples=180] plot(\x,{\x+1/\x});\fill[orange!85!black] (0,0) circle(2.2pt);\end{scope}\node[anchor=west,align=left,text width=5.0cm] at (8.4,4.0) {蓝：$A=\dfrac{x}{x^2+1}$\\[12pt]$D_A=\mathbb R$\\[10pt]紫：$B=x+1/x$\\[10pt]$D_B=\mathbb R\setminus\{0\}$\\[10pt]二者均为奇函数};
\node[text width=12.6cm,align=center] at (7.1,.65) {两函数同坐标尺度；紫线分支不能跨过原式无定义的零点连接。};
\end{tikzpicture}\end{center}
